\documentclass[10pt,a4paper]{amsart}
\usepackage[margin=19mm]{geometry}
\usepackage{amsmath,amssymb,mathtools}
\usepackage[scr=rsfs,cal=euler]{mathalfa}
\usepackage{xfrac}
\usepackage[unicode,hidelinks,bookmarks=false]{hyperref}
\newcommand{\ie}{i.e.\ }
\newcommand{\cf}{cf.\ }
\newcommand{\resp}{respectively\ }
\newcommand{\Subsection}{\subsection}
\providecommand{\nobreakdash}{\nobreak\hbox{-}}
\setlength{\emergencystretch}{2em}
\multlinegap0pt
\newtheorem{comment}{Comment}
\usepackage{cancel}
\usepackage{amssymb}
\usepackage{mathtools}
\usepackage{xfrac}
\newtheorem{prop}{Proposition}
\def\Na{\nabla}
\def\R{\mathbb{R}}
\newcommand{\be}{\begin{equation}}
\newcommand{\ee}{\end{equation}}
\def\h {\hat}
\def\la{\langle}
\def\lf{\left}
\def\ra{\rangle}
\def\ri{\right}
\title{Positive mass theorem for initial data sets with arbitrary ends}
\author{Tin-Yau Tsang}
\date{Source: arXiv:2604.26978v2 (CC BY 4.0)}
\begin{document}
\maketitle

\noindent\emph{Source and licence.} Positive mass theorem for initial data sets with arbitrary ends. Version arXiv:2604.26978v2 (CC BY 4.0), distributed by arXiv under the Creative Commons Attribution 4.0 International licence, http://creativecommons.org/licenses/by/4.0/, as recorded in the arXiv licence metadata for this exact version and shown on the abstract page https://arxiv.org/abs/2604.26978v2. Full licence notice: \texttt{licenses/arxiv-2604.26978-CC-BY-4.0.txt}.\par\smallskip

\noindent\textit{Editorial scope.} Counted unit: the article's prop newDEC (source0.tex lines 495--502). The article's complete original proof of it is delivered with it (source0.tex lines 504--551, one \texttt{proof} environment of 48 lines). No mathematics is written, repaired or re-derived: the statement and the proof are the article's own lines, in the article's own notation, and no retained line is substituted or re-flowed.  No further source-local context is needed: every cross-reference of every retained line resolves inside the two delivered spans.  Nothing was omitted from any retained span, so the package carries no omission record.  No retained line contains a citation, so the wrapper carries no bibliography and no bibliography entry.  No retained line contains a figure environment, an image-inclusion command or drawing code, so no image is delivered and the package declares no asset dependency.  Editorial formatting only, in addition: an AMS \texttt{amsart} preamble that restates the article's own declarations and macros used by the retained lines, namely the environments \texttt{prop} on separate counters, together with the standard packages those lines need.  The retained environments print in this document as Proposition 1 in order of appearance, whereas the article prints the same items with the numbering of its own full document.  The complete native source of the article as distributed in the arXiv e-print for this exact version is bundled unchanged as \texttt{sources/199437/source0.tex}.
\begin{prop}\label{newDEC}
Let $(M,g,k)$ be an initial data set. For a smooth function $h:M\to \R$, define another initial data set $(M,g,\hat k:=k-\frac{h}{n-1}g)$. We have, 
\be
\begin{split}
\h \mu = \mu+\frac{n}{2(n-1)}h^2-h\, tr_g k \,\,\, \text{and} \,\,\, \h J= J+dh. 
\end{split}
\ee 
\end{prop}

\begin{proof}
Let $f\in C^\infty(M)$, consider an initial data set $(M,g,\hat k:=k+fg)$. 
By definition, 
\begin{equation}
\begin{split}
\hat \mu :=& \frac{1}{2}\lf( R_g-|\h k|_g^2+(tr_g \hat k)^2 \ri)\\
%=& \frac{1}{2}\lf( R_g-\la k+fg,k+fg\ra+(tr_g k + nf)^2 \ri)\\
%=& \frac{1}{2}\lf( R_g -|k|_g^2+(tr_gk)^2 -nf^2-2ftr_gk+n^2f^2+2nf tr_g k\ri)\\
=&\mu+\frac{1}{2}n(n-1)f^2+(n-1)ftr_gk.  
\end{split}
\end{equation}

\begin{equation}
\begin{split}
\hat \pi:=&\, \h k-(tr_g \h k)g=\pi+(1-n)fg.\\
%=&k-(tr_g k)g+(1-n)fg\\ 
\end{split}
\end{equation}

\begin{equation}
\begin{split}
\hat J:=& div_{g}(\hat \pi)= J + (1-n) df. \\
%=& div_{g}(\pi)+ (1-n) div_g(fg)\\
%=& J + (1-n) div_g(fg)\\
%=& J + (1-n) g(\Na f, \cdot)\\
\end{split}
\end{equation}



If $f:=\frac{-1}{n-1}h$, then 
\begin{equation}
\begin{split}
\hat \mu 
=&\mu+\frac{1}{2}\frac{n}{(n-1)}h^2-h\, tr_gk.  \\
\hat J= & J+dh. \\
\end{split} 
\end{equation}

\begin{comment}
Therefore, 
\begin{equation}
\begin{split}
\h \mu+\h J(\nu)=&\mu+J(\nu)+\frac{1}{2}\frac{n}{(n-1)}h^2-h\, tr_gk+\Na_{\nu}h. 
\end{split}
\end{equation} 
\end{comment}
\end{proof}

\end{document}
